\section{Rubix：统一运行 Substrate 与 Compliance-as-Code}

Apollo 决定软件应怎样变化，Rubix 提供变化发生的运行地基。这里的 \term{substrate} 指应用之下、具体云或硬件之上的统一运行层。Palantir 当前官方文档把 Rubix 描述为 hardened、autoscaling、highly available 的 Kubernetes implementation，用来承载 AIP、Foundry、Apollo 及客户 workload。它的价值不是隐藏所有差异，而是把必须一致的安全、生命周期和可用性策略固定下来。

\subsection{Kubernetes 统一了什么，仍缺什么}

\term{Kubernetes} 是容器编排系统，负责声明和维持 workload、service、network 与 storage 资源。原生对象可以描述 replicas、volume 和 ingress，却不知道某个国家安全环境的 operator access、节点寿命、合规审批和跨 fleet 发布语义。Rubix 在 Kubernetes 之上建立 opinionated API，将这些组织经验编码为默认行为。

\begin{figure}[H]
\centering
\includegraphics[width=0.94\textwidth]{images/05-rubix-layers.png}
\caption{Rubix 位于应用与异构基础设施之间，Apollo 在其上表达版本与策略。}
\figsource{依据字幕 15:00--20:00 重绘，`images/05-rubix-layers.png`。}
\end{figure}

\begin{knowledgebox}{读图：抽象的边界是“应用意图”}
产品团队声明需要几个 replica、什么持久化、暴露哪个 endpoint、怎样判断 quorum；substrate 再映射成 provider-specific resource、certificate、network policy 和 lifecycle。若应用仍要知道每家云的 ingress、每个机房的 PKI 和每项合规 checklist，抽象就没有真正释放实验速度。
\end{knowledgebox}

\begin{figure}[H]
\centering
\includegraphics[width=0.9\textwidth]{images/official-rubix-k8s.png}
\caption{Palantir 官方图将 Rubix 定位为 hardened、autoscaling、highly available Kubernetes 实现。}
\figsource{Palantir Rubix 官方架构文档，2026-08-11 获取，`images/official-rubix-k8s.png`。}
\end{figure}

\begin{importantbox}{术语消化：Multi-tenancy 与隔离}
Multi-tenancy 指多个客户、团队或 workload 共享底层基础设施，同时保持身份、数据、网络和资源边界。隔离不仅是 namespace；还包括最小权限、network segmentation、secret 管理、resource quota、审计和 noisy-neighbor 控制。共享提高利用率，也扩大配置错误的 blast radius。
\end{importantbox}

\subsection{从人工合规到软件执行}

监管控制通常有合理目标，例如限制 operator、记录变化、加密传输和保持补丁基线。仅靠文档和人工 checklist 时，控制会随环境数量线性增加，也容易产生 drift。\term{Compliance-as-code} 将要求转成 policy、immutable configuration、automated evidence 和阻止违规部署的 gate，让“符合规则”成为系统状态，而不是季度审计时临时收集的文件。

\begin{figure}[H]
\centering
\includegraphics[width=0.94\textwidth]{images/official-rubix-security.png}
\caption{Rubix 官方安全图列出地域控制、审计、zero-downtime upgrade、node cycling、encryption 与 service orchestration。}
\figsource{Palantir Rubix 官方架构文档，2026-08-11 获取，`images/official-rubix-security.png`。}
\end{figure}

\begin{knowledgebox}{读图：安全属性必须落到可观察证据}
“Encryption at rest” 需要能证明 key 与 storage policy；“network monitoring” 需要可查询日志和 alert；“geographic controls” 需要 scheduler 与 data placement 约束。图中的 zero-downtime upgrade 指发布时维持服务，不是分布式训练里的 ZeRO；后者是 Zero Redundancy Optimizer，通过对 optimizer state、gradient 或 parameter 做 sharding 来节省每张 GPU 的显存。安全图是能力目录，不是自动证明。每一项都应关联配置、运行证据、异常处置和 owner。
\end{knowledgebox}

\subsection{Ephemeral node：用替换消除 drift}

\term{Immutable image} 指运行中的 image 不做就地修改，补丁通过构建新 image 并替换 instance 进入系统。\term{Ephemeral node} 是寿命受限、预期被销毁和重建的节点。这样可以减少长期 drift，并迫使 service 具备 failover 能力。Node cycling 还会缩短攻击者保持持久 foothold 的时间，但它不能替代 credential rotation、network policy 和 application security。

\begin{figure}[H]
\centering
\includegraphics[width=0.94\textwidth]{images/06-ephemeral-security.png}
\caption{Ephemerality 把 patching 转成受控 replacement 循环。}
\figsource{依据字幕与当前 Rubix 文档重绘，`images/06-ephemeral-security.png`。}
\end{figure}

\begin{warningbox}{版本差异：72 小时口述与 48 小时当前文档}
课堂字幕中，Sankar 说 container image 随机在约 40--72 小时之间被销毁并从新 image 重建；2026-08-11 获取的 Rubix 官方文档写 nodes cannot live longer than 48 hours。可能是产品策略在不同时间更新，也可能是 image、host 与 node 的口径差异。讲义不把两者混成单一数字，实施时应以目标部署的当前 policy 为准。
\end{warningbox}

\begin{table}[H]
\centering
\begin{tabular}{L{0.19\textwidth}L{0.32\textwidth}L{0.39\textwidth}}
\toprule
机制 & 获得的能力 & 新增要求 \\
\midrule
Immutable image & 基线可复现、减少就地 drift & 构建与签名链可信 \\
Bounded lifetime & 限制持久化、强制更新 & Workload 可重调度、状态外置 \\
Policy-driven drain & 减少同时驱逐造成的抖动 & 拓扑、容量和任务优先级可见 \\
Blue/green rollout & 并行验证、zero-downtime 切换 & 双份容量与健康判定 \\
Audit logging & 追溯 operator 与系统行为 & 日志完整性和访问控制 \\
\bottomrule
\end{tabular}
\caption{安全自动化会把隐含人工依赖转成新的软件与容量要求。}
\end{table}

\subsection{本章小结}

Rubix 用统一 Kubernetes substrate 承担安全、隔离、生命周期和跨环境一致性，Apollo 再利用这个 substrate 执行版本与配置变化。Compliance-as-code 和 ephemeral infrastructure 降低人工 toil，却要求 workload 真正支持失效、重建和外置状态。抽象不是消除复杂度，而是把复杂度交给能集中验证的平台团队。

